\begin{afterword}
\summaryhead

The essay asks why people who call themselves rationalists are not visibly more capable than anyone else. It starts from a common objection, that rational people do not seem happier. The author replies that the people the objectors have in mind have little rationality. Even the best researchers in the psychology of reasoning fall short of what should be possible.

His explanation is that rationality has never been systematized, trained and tested. A mathematician, an engineer or an athlete draws on organized knowledge and long practice. Rationalists improvise, and the author includes himself: he is strong in one skill he built alone and weak in others.

The model is the martial arts. Schools of fighting exist because people wanted to become formidable, and they improved because they could tell when a blow landed. Rationality lacks all three conditions: people who work together, ways to verify training, and, first, the sense that more is possible.

\responsehead

The essay is an open call for rationality schools, modelled on martial-arts dojos, with training programs tested against each other. Its diagnosis is that the art does not exist because no one has systematized it, tested it or wanted it enough. Each part of that diagnosis has a weaker spot than the essay admits.

It opens with an objection: rational people do not seem happier. The answer is that the people the objector has in mind are not very rational. The essay never says how rational would be enough, and the bar rises as it goes, past anyone who can derive Bayes's Theorem and past the most senior researchers the author admires. Whatever example an objector brings, the answer would be the same. The objection is set aside without being tested, and the essay's question changes from happiness to formidability.

The diagnosis assumes that there is a large general skill of thinking waiting to be trained. The essay does not consider the alternative: that beyond intelligence, knowledge and practice in particular fields, there may be little general skill to train, and so little for a school to add. The evidence the essay offers, that no one has yet become visibly formidable through rationality, fits that alternative at least as well as it fits the essay's own explanation.

Its account of the research is thin at the point where it matters. It says that debiasing studies amount to an hour of practice checked two weeks later, and cites none. Nisbett and his colleagues had studied what years of university training do to reasoning and concluded that ``reasoning is taught and generalizable''; their short sessions also had measured effects. None of this was the randomized comparison of programs the essay wants. But the question it says nobody asks had been asked, with results.

The martial-arts analogy carries the argument, and it rests on a claim that is not true of martial arts in general: that the schools improved because they competed in realistic contests. Some schools reject competition as contrary to their spirit, and flourish anyway. The analogy the essay chose shows how a discipline can grow without the testing the essay calls essential.

In short: a call for rationality schools that sets aside its opening objection by raising the bar, never considers that there may be less general skill to train than it hopes, overlooks the studies that had measured training, and leans on an analogy that does not hold.
\end{afterword}
